\subsection{Усреднение функций, имеющих обобщённые производные}

\begin{lemma}[Л1]\label{lem:avg-commute}
Пусть $u \in \Lp{1}(\Om)$ и $\exists\, D^\alpha u \in \Lp{1}(\Om)$. Тогда
\[
  (D^\alpha u)_\rho(x) = \big(D^\alpha u_\rho\big)(x) ,
\]
где $x \in \Om_\rho$, $\dist(\dOm, x) > 0$, а $\rho$ таково, что
$0 < \rho < \dist(\dOm, x)$.
\end{lemma}

\begin{idea}
Усреднение и обобщённое дифференцирование коммутируют — но только вдали от границы:
шар $\ball{x}{\rho}$ должен целиком лежать в $\Om$, иначе ядро «высунется» наружу,
где $u$ не определена.
\end{idea}

\begin{proof}
Проводим выкладку для $\partial/\partial x_i$; общий случай аналогичен.

При фиксированном $x$ функция $y \mapsto \omega_\rho(x-y)$ обращается в нуль вне
$\ball{x}{\rho}$, то есть финитна в $\Om$. Дифференцируя под знаком интеграла
(свойство 3 усреднения),
\[
  \dd{u_\rho}{x_i}(x)
  = \dd{}{x_i}\left( \int\limits_\Om \omega_\rho(x-y)\, u(y)\, dy \right)
  = \int\limits_\Om \dd{}{x_i}\,\omega_\rho(x-y)\, u(y)\, dy .
  \tag{$\ast\ast$}\label{eq:dstar}
\]
Далее, ядро зависит от разности, поэтому
\[
  \dd{}{x_i}\big(\omega_\rho(x-y)\big) = -\,\dd{}{y_i}\big(\omega_\rho(x-y)\big) ,
\]
и, подставляя это в \eqref{eq:dstar} и применяя определение обобщённой производной
к пробной функции $\varphi \equiv \omega_\rho(x - \cdot)$,
\[
  \eqref{eq:dstar} = -\int\limits_\Om \dd{\omega_\rho(x-y)}{y_i}\, u(y)\, dy
  = \int\limits_\Om \omega_\rho(x-y)\, \dd{u}{y_i}(y)\, dy
  = \left( \dd{u}{x_i} \right)_{\!\rho}(x) .
\]

Для произвольного мультииндекса всё то же самое, только знак набегает $|\alpha|$ раз:
\[
  \boxed{\;D^\alpha_x\, \omega_\rho(x-y) = (-1)^{|\alpha|}\, D^\alpha_y\, \omega_\rho(x-y)\;}
\]
\end{proof}

\begin{lemma}[Л2]\label{lem:avg-conv-loc}
Пусть $1 \le p < \infty$, $u \in \Lp{p}(\Om)$ и $\exists\, D^\alpha u \in \Lp{p}(\Om)$.
Тогда для любого $\Om' \Subset \Om$
\[
  u_\rho \to u \ \text{ в } \Lp{p}(\Om'), \qquad
  D^\alpha u_\rho \to D^\alpha u \ \text{ в } \Lp{p}(\Om') .
\]
\end{lemma}

\begin{warn}
В тетради во второй строке заключения сначала было написано $\Lp{p}(\Om)$ — зачёркнуто
и исправлено на $\Lp{p}(\Om')$. Сходимость именно локальная: около границы лемма~\ref{lem:avg-commute}
неприменима.
\end{warn}

\begin{proof}
Знаем, что $u_\rho \to u$ в $\Lp{p}(\Om)$ (теорема~\ref{thm:conv}), а значит и
$u_\rho \to u$ в $\Lp{p}(\Om')$.

Так же $(D^\alpha u)_\rho \to D^\alpha u$ в $\Lp{p}(\Om')$ — это та же теорема,
применённая к функции $D^\alpha u \in \Lp{p}(\Om)$.

Пусть $\dist(\Om', \dOm) = \rho_0 > 0$ и $0 < \rho < \rho_0$. Тогда по
лемме~\ref{lem:avg-commute} в $\Om'$ выполнено $(D^\alpha u)_\rho = D^\alpha u_\rho$,
откуда
\[
  D^\alpha u_\rho \to D^\alpha u \ \text{ в } \Lp{p}(\Om') . \qedhere
\]
\end{proof}

\begin{lemma}[Л3]\label{lem:const}
Пусть $u \in \Lloc(\Om)$, область $\Om$ не обязательно ограничена, и $Du = 0$
(все обобщённые производные первого порядка равны нулю). Тогда $u = \mathrm{const}$
почти всюду в $\Om$.
\end{lemma}

\begin{proof}
\textbf{1) Случай $u \in \Lp{1}(\Om)$, $\Om$ ограничена.}

Для $\eps > 0$ положим $\Om_\eps = \{x \in \Om \mid \dist(x, \dOm) > \eps\}$.
По лемме~\ref{lem:avg-commute} при $0 < \rho < \eps$ в $\Om_\eps$
\[
  \left( \dd{u}{x_i} \right)_{\!\rho} = \dd{u_\rho}{x_i} ,
  \qquad\text{причём}\qquad
  \left( \dd{u}{x_i} \right)_{\!\rho}
  = \int\limits_{\Om_\eps} \omega_\rho(x-y)\, \dd{u}{x_i}(y)\, dy = 0 ,
\]
так как $\partial u/\partial x_i = 0$. Значит, все частные производные гладкой функции
$u_\rho$ в $\Om_\eps$ нулевые, и
\[
  u_\rho(x) = C(\rho) \quad \text{в } \Om_\eps ,
\]
поскольку $u_\rho$ — гладкая функция.

\emph{Покажем, что $C(\rho) < \infty$.} По теореме~\ref{thm:norm}
\[
  |C(\rho)|\cdot \mu(\Om_\eps) = \int\limits_{\Om_\eps} |u_\rho|\, dx
  \le \int\limits_\Om |u_\rho|\, dx \le \int\limits_\Om |u|\, dx < \infty ,
\]
откуда
\[
  |C(\rho)| \le \frac{1}{\mu(\Om_\eps)} \int\limits_\Om |u|\, dx < \infty .
\]
Пусть $C(\rho_k) \to C_0$ при $\rho \to 0$ (выбрали сходящуюся подпоследовательность).
Тогда
\begin{multline*}
  \int\limits_{\Om_\eps} |u - C_0|\, dx
  = \int\limits_{\Om_\eps} |u - C(\rho) + C(\rho) - C_0|\, dx \\
  \le \int\limits_{\Om_\eps} |u - u_\rho|\, dx + \int\limits_{\Om_\eps} |C(\rho) - C_0|\, dx
  \le \underbrace{\int\limits_{\Om_\eps} |u - u_\rho|\, dx}_{\to\, 0}
    + \underbrace{|C(\rho) - C_0|\cdot \mu(\Om_\eps)}_{\to\, 0} \longrightarrow 0 ,
\end{multline*}
то есть $\int_{\Om_\eps} |u - C_0|\, dx = 0$, а значит $u = C_0$ п.\,в. в $\Om_\eps$.

Положим $\eps = \tfrac1k$, то есть $\Om_\eps = \{x \in \Om \mid \dist(x,\dOm) > \tfrac1k\}$.
Тогда $\mu\big(\Om_{1/k}[u \neq C_0]\big) = 0$ и
\[
  \Om[u \neq C_0] = \bigcup_{k \in \N} \Om_{1/k}[u \neq C_0]
  \quad\Longrightarrow\quad u = C_0 \ \text{п.\,в. на } \Om .
\]

\textbf{2) Случай $u \in \Lloc(\Om)$, $\Om$ не обязательно ограничена.}

Для любого $\Om' \Subset \Om$ имеем $u \in \Lp{1}(\Om')$, а для этого случая всё уже
доказано, так что $u = \mathrm{const}$ в $\Om'$. Возьмём
\[
  \Om_k = \Big\{x \in \Om \ \Big|\ \dist(x, \dOm) > \tfrac1k \Big\} \cap \ball{0}{k},
  \qquad \Om_k \Subset \Om, \qquad \Om = \bigcup_{k \in \N} \Om_k .
\]
Тогда $\mu\big(\Om_k[u \neq C_0]\big) = 0$, а
\[
  \Om[u \neq C_0] = \bigcup_{k=1}^{\infty} \Om_k[u \neq C_0]
  \quad\Longrightarrow\quad \mu\big(\Om[u \neq C_0]\big) = 0 . \qedhere
\]
\end{proof}

\begin{idea}
Пересечение с шаром $\ball{0}{k}$ во втором пункте нужно именно потому, что $\Om$ может
быть неограниченной: одно условие $\dist(x,\dOm) > 1/k$ компактности не даёт.
\end{idea}

\begin{theorem}[замена переменных]\label{thm:change-vars}
Пусть $\Om, \Om' \subset \R^n$, а $\varphi : \Om' \to \Om$ — дифференцируемая биекция
с непрерывными первыми производными и ненулевым якобианом. Пусть
\[
  u \in \Lloc(\Om), \qquad \exists\, \dd{u}{x_i} \in \Lloc(\Om), \quad i = \overline{1,n} .
\]
Тогда $v(y) = u(\varphi(y)) \in \Lloc(\Om')$ и $v$ имеет обобщённые производные
из $\Lloc(\Om')$:
\[
  \dd{v}{y_i} = \sum_{k=1}^{n} \dd{u}{x_k}\big(\varphi(y)\big)\, \dd{\varphi_k(y)}{y_i} .
\]
\end{theorem}

\begin{proof}
\plusplus
\end{proof}
